Der menschliche Körper ist kein passiver Gegenstand, der einfach Wärme austauscht, bis die Temperaturen gleich sind:
\begin{itemize}
\item Durch den Stoffwechsel erzeugt er dauernd Wärme (in Ruhe etwa \PKO), und die Temperaturregulation hält die Kerntemperatur bei etwa \TKO. Ist das Wasser kühler, ersetzt der Körper die abgegebene Wärme laufend (und drosselt die Durchblutung der Haut).
\item Das Badewasser hat eine viel grössere Masse als der Körper. Die Wärme, die der Körper abgibt oder aufnimmt, ändert die Wassertemperatur kaum, und das Wasser gibt gleichzeitig Wärme an die Luft und die Wanne ab.
\item Der Wärmeaustausch ist langsam: Bis zu einem vollständigen Temperaturausgleich würde es viel länger dauern als ein Bad.
\end{itemize}
Nur die Haut nimmt annähernd die Wassertemperatur an; im Körperinnern bleibt die Temperatur nahezu konstant.
